\section{本章小结}
\label{sec:hj-summary}

本章围绕一个问题展开：当方程只限制局部的梯度时，场值由什么决定。到达时间给出了具体的图景，波前法向推进的速度决定$T$梯度的长度；令$u=T-t$，到达时间问题就进入时间依赖的方程，Hamilton量$H$说明斜率$\bm{p}$的地方高度下降多快。两类问题的数据不同，边值问题还要求数据相容。

光滑时，方程化为特征方程组。位置和斜率服从Hamilton方程，值靠一次积分带上，这样可以局部重建经典解。对波前而言，特征线就是射线，$u$沿线不变。

特征线会相交。一维中曲率按式~\eqref{eq:hj-riccati}发散，经典解在有限时刻结束，而斜率和场值依然有界。此后同一位置有多个候选，几乎处处满足方程的解不唯一，凹槽就是例子。

黏性解在接触点处提出不等式：从上方贴得到$\leq0$，从下方贴得到$\geq0$。它不做积分，也不移走导数，而且对方程的符号敏感。凹折角通过检验，凸折角不通过，这和“只有第一次到达才算数”是一致的。

比较原理把这些局部的不等式变成整体结论：次解和上解的次序从数据传到全部时刻，因此黏性解唯一，且对初值的依赖在上确界范数下不扩张。初值问题要有增长和连续性条件，边值问题要有$n>0$和相容的数据，两种情形的严格性来源不同，不能互相搬用。

存在性有两条途径。对凸且超线性增长的$H$，Hopf--Lax公式写出解，其极小点就是到达的特征线，取最小值就是选下包络，多个极小点的地方就是不可微点。消失黏性把解看成小扩散方程的极限，$H=\lVert\bm{p}\rVert_2^2/2$时Cole--Hopf变换让这件事可以精确计算；一般情形下依靠一致估计、局部一致收敛下的稳定性和比较原理。有限黏性下解仍有误差，加细网格也不能消除它。

表~\ref{tab:hj-summary}汇总各工具解决的问题和适用条件。

\begin{BookInlineTable}
  \centering
  \begin{tabular}{lll}
    \toprule
    工具 & 解决的问题 & 条件\\
    \midrule
    特征方程组 & 光滑范围内重建解 & $H$、$u_0$光滑，局部\\
    黏性解 & 不可微点上的选择 & 连续函数\\
    比较原理 & 唯一性与稳定性 & 增长、连续性、相容性\\
    Hopf--Lax公式 & 显式表示与存在性 & $H$凸且超线性\\
    消失黏性 & 光滑近似与解的选择 & 一致估计\\
    \bottomrule
  \end{tabular}
  \caption{本章各工具的作用。}
  \label{tab:hj-summary}
\end{BookInlineTable}

本章只做方程分析，所有结论都是有限时间内的。这里的$H$是给定的，Hopf--Lax公式中的极小化也只是解的表示。在决策问题中，Hamilton量由系统和代价确定，价值函数满足Hamilton--Jacobi--Bellman方程，需要进一步讨论价值函数与策略的验证，可参考\BookChapterRef{chap:continuous-time-optimal-control}{卷十第6章“连续时间最优控制”}。
